\chapter{The capstone theorem and the green rungs}
\label{chap:capstone}

Every earlier chapter speaks about \emph{relations}: erasability and $\Erases$
(Chapter~\ref{chap:erases}), the forward simulation for $\Erases$ (Chapter~\ref{chap:correctness}),
the Lean-specific lowering $\Lower$ and the emitted program's well-formedness
(Chapter~\ref{chap:lowering}), the shipping eraser's specification and its run model
(Chapters~\ref{chap:eraser} and~\ref{chap:run}), the refinement of \code{visitExpr}
(Chapter~\ref{chap:refinement}), and the bridge that turns a run into an erasure
(Chapter~\ref{chap:coldstart}). This chapter states the theorem those pieces were built for and then
instantiates it eight times at runs of the real \code{\#erase} command. The theorem is
\code{shipping\_erase\_correct\_firstorder} (\code{LeanToLambdaBox/Capstone.lean}): under fifteen
named hypotheses it says that the program \code{Erasure.erase} emits is the $\Lower$ image of a
specification environment that erases the prepared source term, that it satisfies what peregrine's
first pass reads of a program, and that every \emph{first-order} answer the source evaluation
produces is reproduced by the emitted program under \lbox{}'s own call-by-value semantics, uniquely
and free of $\Box$.

The chapter's second half is the \emph{green ladder} (\code{LeanToLambdaBox/Green.lean}, 1552
lines): eight theorems \code{green\_G1} through \code{green\_G8}, each an instance of the capstone
at a committed \code{\#erase} run whose emitted environment and term are transcribed into the file
and whose conclusion ends in a \textbf{literal} \lbox{} peano numeral. The literal answer is an
anti-vacuity device: a correctness statement of this shape is cheap to satisfy if the target answer
may be $\Box$ or a stuck term, and pinning it to a numeral with the certified evaluator and then
forcing the theorem's existential with determinism of evaluation rules both out. The ladder is
nevertheless conditional --- the Lean environment \code{lenv} and the model environment \code{env}
are universally quantified at every rung, and the run equation \code{hrun} is an equation about an
opaque \code{IO.RealWorld} token that no Lean term can inhabit --- so what a rung delivers is:
\emph{these hypotheses are consistent with a literal answer, and every hypothesis a computation can
settle is settled by one.}

Between the two halves sits the \emph{reification} machinery of
\code{LeanToLambdaBox/Witness/SourceTable.lean} and \code{LeanToLambdaBox/Witness/TrWitness.lean}.
The verification never reads a \code{Lean.Environment} and never ingests a committed artefact: the
term elaborator \code{reify\%} copies the slice of the environment a rung needs into a literal
source table, so that the rung's side conditions are kernel computations (\code{by rfl}, \code{by
decide +kernel}) rather than assumptions about opaque data. The faithfulness of that copy is one
named hypothesis, table adequacy (the capstone's \code{htbl}), mechanised outside the kernel by
\code{lake exe reify --check}. \code{TrWitness.lean} then turns a tabled monomorphic constant into
the model translation the capstone's \code{hwt} asks for, so that binder too costs nothing at a
rung. The chapter must be read together with Chapter~\ref{chap:trust}: the capstone is a composition
with exactly one hole, its binder \code{hbridge}, which packages the two environment-level results
($\ErasesEnv$ and $\LowerEnv$, the two fields of \code{ErasureBridge}) that no theorem in the
repository produces for a run of the shipping eraser, and which is inhabited at \textbf{no} rung.

\section{What the capstone composes}
\label{sec:capstone:overview}

Read as a composition, the capstone spends four ingredients. \emph{The run split}
(\code{erase\_run\_ok}, Chapter~\ref{chap:coldstart}) factors one shipping run into a preparation
and a \code{visitExpr} run from the empty state. \emph{The bridge} (\code{erasure\_bridge\_of\_run},
Chapter~\ref{chap:coldstart}) turns that \code{visitExpr} run into an intermediate \lbox{} term
$t_0$ with an $\Erases$ derivation and a $\Lower$ derivation. \emph{The simulation}
(\code{erases\_correct}, Chapter~\ref{chap:correctness}) is applied once, at the applied spine, and
yields the target evaluation. \emph{The first-order core} (\code{firstorder\_erases\_core} and
\code{noBox\_lower\_of\_foSpine}, Chapter~\ref{chap:correctness}) gives the answer's
constructor-spine shape, the uniqueness of its erasure and the box-freedom of its lowering. What the
composition does \emph{not} supply is the pair of environment-level facts about the shipping
eraser's own registration behaviour; those enter as the binder \code{hbridge}.

The capstone's measured axiom footprint is 33 names (the committed fixture
\code{test/ledger.expected}), identical to every rung's: four standard axioms (\code{propext},
\code{sorryAx}, \code{Classical.choice}, \code{Quot.sound}) and 29 non-standard ones, of which two
are \code{bv\_decide} certificates and two are lean4lean pointer-equality axioms, the rest
reflection axioms for \code{Lean.Expr}, \code{Lean.Level}, the persistent containers and syntax. The
\code{sorryAx} is inherited from the pinned lean4lean checkout; Chapter~\ref{chap:trust} states what
it bounds.

\section{The reified source table}
\label{sec:capstone:table}

\begin{definition}[The reified source table]\label{def:Witness-SourceTable}
\lean{LeanToLambdaBox.Witness.SourceTable, LeanToLambdaBox.Witness.ReifiedDecl, LeanToLambdaBox.Witness.ReifiedInduct, LeanToLambdaBox.Witness.ReifiedCtor, LeanToLambdaBox.Witness.SourceTable.decl?, LeanToLambdaBox.Witness.SourceTable.ind?, LeanToLambdaBox.Witness.SourceTable.body?} \leanok
A \emph{source table} is the reified slice of a Lean elaboration environment that the verification
reads in the environment's place: a list of named \emph{reified declarations} and a list of named
\emph{reified inductive types}. A reified declaration records the level parameters and the type the
kernel constant carries, together with an optional body; a reified inductive type records its level
parameters, its type, its parameter and index counts, the names of its mutual block, and for each
constructor its name, type, position and parameter/field split. The body column is optional because
the eraser only ever opens the body of a definition or an opaque constant: axioms, theorems,
quotient primitives, constructors, recursors and \code{casesOn}-like constants it eliminates with
rather than unfolds, and those are tabled with no body. The lookup interface is the three Lean
functions \code{SourceTable.decl?}, \code{SourceTable.ind?} and the partial map
\code{SourceTable.body?}, which is the $\delta$-column that the source evaluation, $\ErasesEnv$,
\code{CompilerBodies} and \code{SpecEnv} all read.
\end{definition}

\begin{remark}
The table has \textbf{no} relevance-oracle column (relevance is discharged, not tabled) and no
configuration column: it is built under one fixed configuration, and a rung states its configuration
obligation separately. The source evaluation reads this table's compiler-body column, not the
kernel's $\delta$ rules --- design tag N8.
\end{remark}

\begin{definition}[The compiler's view of a declaration]\label{def:Witness-compilerInfo}
\lean{LeanToLambdaBox.Witness.compilerInfo?, LeanToLambdaBox.Witness.compilerValue?, LeanToLambdaBox.Witness.fixBlock?, LeanToLambdaBox.Witness.reifiesBody} \leanok
\code{compilerInfo? lenv n} is the constant information the \emph{code generator} reads for
\code{n}: the \code{\_unsafe\_rec} companion the elaborator emits for a recursive definition when
there is one, and \code{n} itself otherwise. \code{compilerValue?} takes that constant's value with
\code{allowOpaque := true}, exactly as the eraser's mutual-block entry point does; \code{none} is
the case in which the eraser emits an axiom instead. \code{fixBlock? lenv n} is the mutual block for
which the eraser installs a fix-variable map at \code{n} --- \code{none} when it takes its
non-recursive exit or emits an axiom, and otherwise the block's members after the
\code{\_unsafe\_rec} stripping the run applies. \code{reifiesBody} is the Boolean test of a
\code{ConstantInfo}'s kind that this same dispatch uses to decide whether a body should be sought
at all: true exactly at \code{Lean.ConstantInfo.defnInfo} and \code{Lean.ConstantInfo.opaqueInfo},
false at every other constructor (axioms, theorems, quotient primitives, constructors, recursors),
none of which the eraser reads a body of.
\end{definition}

\begin{remark}
This is the formal locus of the \emph{compiler-versus-kernel bodies} scope restriction: the eraser
never reads the kernel body of a structurally recursive definition, which eliminates with
\code{brecOn}, but the compiler body, which calls itself directly. \code{fixBlock?} is what
\code{lake exe reify --blocks} computes in order to mechanise two of the three clauses of
\code{TableBlocks}.
\end{remark}

\begin{lemma}[The compiler's view is Lean's own]\label{lem:Witness-compilerInfo-eq}
\lean{LeanToLambdaBox.Witness.compilerInfo?_eq} \leanok
\uses{def:Witness-compilerInfo}
For every name \code{n}, Lean's own \code{Compiler.LCNF.getDeclInfo? n} is the monadic action that
reads the ambient environment and returns \code{compilerInfo?} of it at \code{n}.
\end{lemma}

\begin{proof} \leanok
\uses{def:Witness-compilerInfo}
Definitional: the two sides are the same unfolding of \code{getDeclInfo?}, and the Lean proof is
\code{rfl}. This is the only link between the table's model of the code generator's view and the
primitive the shipping eraser actually calls.
\end{proof}

\begin{definition}[Alpha-equivalence of Lean expressions]\label{def:Witness-Expr-AlphaEq}
\lean{LeanToLambdaBox.Witness.Expr.AlphaEq} \leanok
\code{Expr.AlphaEq} is structural equality of \code{Lean.Expr} that ignores binder names and binder
info, and nothing else: twelve congruence arms over \code{bvar}, \code{fvar}, \code{mvar},
\code{sort}, \code{const}, \code{lit}, \code{app}, \code{lam}, \code{forallE}, \code{letE},
\code{proj} and \code{mdata}, with the three binder arms quantifying the two binder names and the
two binder annotations away. It deliberately has \textbf{no} metadata-blind arm: \code{.mdata d e}
and \code{e} stay unrelated, because the source semantics has no metadata rule and a relation
identifying them would falsify an alpha-transport of that semantics.
\end{definition}

\begin{remark}
Extensionally this is \code{Lean.Expr.eqv}, which is \code{@[extern]} and opaque, so \textbf{no
theorem connects the two}. \code{Expr.alphaEqB} is the arm-for-arm Boolean counterpart (a
\code{partial def}), and \code{lake exe reify --check} cross-checks it against \code{Lean.Expr.eqv}
on every compared pair, reporting a disagreement as a mismatch. Chapter~\ref{chap:trust} carries the
gap as a finding.
\end{remark}

\begin{definition}[Pinning a tabled constant and a tabled inductive type]\label{def:Witness-ReifiedDecl-Pinned}
\lean{LeanToLambdaBox.Witness.ReifiedDecl.Pinned, LeanToLambdaBox.Witness.ReifiedInduct.Pinned} \leanok
\uses{def:Witness-SourceTable}
A tabled constant is \emph{pinned} in an environment when the environment knows that name with
exactly the level parameters and the type the table records. A tabled inductive type is pinned when
the environment knows the name as an inductive of that name, listed in its own mutual block, with
the recorded level parameters, type, parameter count, index count and block; and when each tabled
constructor is the environment's constructor of that type at its position in the tabled list, with
the recorded type, constructor index (equal to its position), parameter count (equal to the block's)
and field count.
\end{definition}

\begin{remark}
The positional clauses of the inductive pin are the kernel's own indexing invariants restated; they
are what carries the table's arithmetic --- constructor arities and indices --- to the specification
bundle's kernel-field clause.
\end{remark}

\begin{definition}[The body column, pinned up to alpha]\label{def:Witness-ReifiedDecl-Prepared}
\lean{LeanToLambdaBox.Witness.ReifiedDecl.Prepared} \leanok
\uses{def:Witness-compilerInfo, def:Witness-Expr-AlphaEq, def:Erasure-prepare_erasure}
A tabled declaration is \emph{prepared} in an environment when, for every body $b$ the table records
for it, the code generator reads a value $v$ for that name in the environment and \emph{every}
successful run of \code{Erasure.prepare\_erasure} on $v$, under a configuration with \code{csimp}
switched off, returns a term alpha-equivalent to $b$. This is the one clause of table adequacy that
is a statement about monadic runs rather than about the output of an environment lookup.
\end{definition}

\begin{remark}
Stated with equality in place of alpha-equivalence the clause would be \textbf{uninhabited} wherever
Lean's \code{inlineMatchers} pass fires --- \code{Nat.add} included --- because the inliner names
its inlined \code{let} binders from the name generator. The price is that a consumer of a tabled
body owes an alpha-transport, and $\Lower$ is not alpha-closed on its source; this is one of the
three measured obstructions to closing \code{hbridge}. Measured consequence: \code{lake exe reify
--check} passes on all eight rung tables and reports five bodies matching only up to binder names at
G7 and G8 (\code{Nat.add}, \code{Nat.mul}, \code{Nat.pow}, \code{Nat.pred}, \code{Nat.sub}).
\end{remark}

\begin{definition}[Table adequacy]\label{def:Witness-SourceTableAdequate}
\lean{LeanToLambdaBox.Witness.SourceTableAdequate, LeanToLambdaBox.Witness.SourceTableAdequate.decls, LeanToLambdaBox.Witness.SourceTableAdequate.inds} \leanok
\uses{def:Witness-ReifiedDecl-Pinned, def:Witness-ReifiedDecl-Prepared, def:Witness-SourceTable}
A table is \emph{adequate} for an environment when every tabled constant is pinned in it and
prepared, and every tabled inductive type is pinned in it. Adequacy is the single interface between
the live elaboration environment and everything downstream. It is not decidable, because no Lean
term denotes the ambient environment, so it is a permanent named hypothesis of every theorem that
reads a table: the capstone's binder \code{htbl}, for which Chapter~\ref{chap:trust} carries the
assumption node.
\end{definition}

\begin{lemma}[A tabled body is a prepared compiler body]\label{lem:Witness-SourceTableAdequate-body-prepared}
\lean{LeanToLambdaBox.Witness.SourceTableAdequate.body?_prepared} \leanok
\uses{def:Witness-SourceTableAdequate, def:Witness-SourceTable}
If a table is adequate for an environment and its body column answers $b$ at a name, then the code
generator reads a value for that name in the environment whose every successful \code{csimp}-off
preparation run returns a term alpha-equivalent to $b$.
\end{lemma}

\begin{proof} \leanok
\uses{def:Witness-SourceTableAdequate, def:Witness-ReifiedDecl-Prepared}
Case on the association-list lookup behind \code{SourceTable.body?}. In the \code{none} case the
hypothesis is the impossible equation \code{none = some b}. In the \code{some d} case, a successful
lookup exhibits a member of the declaration column, so the \code{decls} field of adequacy applies at
that member, and its second component is the prepared clause, read at $b$. This is the form in which
every consumer of a tabled body uses adequacy.
\end{proof}

\begin{definition}[Reification: the \code{reify\%} elaborator]\label{def:Witness-elabReify}
\lean{LeanToLambdaBox.Witness.elabReify, LeanToLambdaBox.Witness.Reify.table, LeanToLambdaBox.Witness.Reify.visit, LeanToLambdaBox.Witness.Reify.State, LeanToLambdaBox.Witness.reifyConfig} \leanok
\uses{def:Witness-SourceTable, def:Witness-compilerInfo, def:Erasure-prepare_erasure, def:Erasure-ErasureConfig}
The term syntax \code{reify\% c1, ..., cn} elaborates to the source table of the dependency closure
of the named constants, read out of the environment of the elaboration it runs in. The worklist
\code{Reify.visit} marks a name seen and dispatches on its constant kind: an inductive type is
reified together with its whole mutual block, with recursion into the constructor types; a
constructor is recorded without a body and its type traversed; a recursor is recorded without a body
and its block members visited; any other constant has its compiler value run through
\code{Erasure.prepare\_erasure} under \code{reifyConfig} and the prepared body recorded, with
recursion into the constants that body names. \code{Reify.table} runs the worklist over a list of
roots and sorts both columns. \code{reifyConfig} is the configuration those preparation runs use: it
sets \code{csimp := false}, the one field preparation reads, and every correctness statement needs
it off, since \code{csimp} replaces bodies by tail-recursive variants that are not the declaration's
own.
\end{definition}

\begin{remark}
\code{reify\%} is the mechanism that makes a rung's side conditions \emph{kernel computations on a
literal} rather than assumptions about an opaque artefact; this is the chapter's main methodological
point. \code{Reify.visit} is a \code{partial def} --- termination by finiteness of the environment
is argued in its doc comment, not proved --- but no proposition reduces it, and its output is only
ever a literal whose faithfulness is table adequacy. The \code{syntax} declaration behind
\code{reify\%} is not a citable Lean declaration, so only the elaborator and the definitions appear
above.
\end{remark}

\begin{definition}[The external table check]\label{def:Witness-SourceTable-check}
\lean{LeanToLambdaBox.Witness.SourceTable.check, LeanToLambdaBox.Witness.checkDecl, LeanToLambdaBox.Witness.checkInd, LeanToLambdaBox.Witness.TableMismatch, LeanToLambdaBox.Witness.TableNote, LeanToLambdaBox.Witness.Expr.alphaEqB} \leanok
\uses{def:Witness-SourceTableAdequate, def:Witness-Expr-AlphaEq, def:Witness-elabReify}
\code{SourceTable.check} is the executable comparison of a table against the live environment:
\code{checkDecl} compares level parameters, compares the type with the structural
\code{Lean.Expr.equal}, and compares the body against a fresh preparation run of the compiler value;
\code{checkInd} compares block data and each constructor positionally. Disagreements are reported as
one of twelve \code{TableMismatch} cases, and the one agreement worth printing --- a body matching
only up to binder names --- as the single \code{TableNote} case.
\end{definition}

\begin{remark}
The check mechanises strictly \textbf{less} than adequacy asserts, as its own doc comment says: the
body column is re-derived by running preparation \emph{once}, whereas the prepared clause quantifies
over all runs, and no theorem connects \code{Lean.Expr.eqv} (used for the body comparison) to
\code{Expr.AlphaEq}. It is therefore wrong to say that \code{lake exe reify --check} \emph{decides}
table adequacy. The module also carries a kernel self-test, including a deliberately stale table the
external check must reject; those examples are anonymous and cannot be cited.
\end{remark}

\section{Translation witnesses: discharging \code{hwt}}
\label{sec:capstone:trwitness}

\begin{lemma}[A declared constant translates]\label{lem:Witness-trExprS_const}
\lean{LeanToLambdaBox.Witness.trExprS_const, LeanToLambdaBox.Witness.trExprS_const_nil} \leanok
\uses{def:ErasureSpec, def:lean4lean-grounding}
Under the specification bundle, a constant the ambient environment declares at a safety the bundle
accepts translates, at any modelled local context, to the model constant of the same name; the level
arguments are translated pointwise and the arity condition is read at the declaration's own level
parameters. The monomorphic instance: a constant with no level parameters, applied to no level
arguments, is its own translation.
\end{lemma}

\begin{proof} \leanok
\uses{def:ErasureSpec, asm:ErasureSpec-decl_adequate, def:lean4lean-grounding}
The bundle's \code{decl\_adequate} field turns the environment lookup into a model constant together
with the equality of universe-variable counts and the translation of the declared type; the constant
arm of lean4lean's translation relation then applies with the translated level list and the arity
equation. The monomorphic instance instantiates the level list with the empty list. This lemma is
the base of the \emph{table route}, the only way a rung obtains the capstone's \code{hwt} without a
binder.
\end{proof}

\begin{lemma}[The declared type of a monomorphic constant, in the model]\label{lem:Witness-hasType_const_nil}
\lean{LeanToLambdaBox.Witness.hasType_const_nil} \leanok
\uses{lem:Witness-trExprS_const, def:ErasureSpec, def:lean4lean-grounding}
If the environment declares a constant with no level parameters whose declared type is itself a
constant, then the model types that constant at the same constant.
\end{lemma}

\begin{proof} \leanok
\uses{def:ErasureSpec, asm:ErasureSpec-decl_adequate, def:lean4lean-grounding}
The \code{decl\_adequate} field gives the model constant and the translation of its declared type;
inverting that translation at a constant type identifies the model type with the model constant of
the same name and forces the level list to be empty, after which the constant rule of the model
typing judgement applies with the universe-count equation read off the declaration.
\end{proof}

\begin{lemma}[A tabled constant is pinned]\label{lem:Witness-tableDeclPin}
\lean{LeanToLambdaBox.Witness.tableDeclPin} \leanok
\uses{def:Witness-SourceTableAdequate, def:TableSafe}
If a table is adequate for an environment and the environment's tabled declarations are safe, then a
successful table lookup yields a constant the environment knows, at an accepted safety, with the
level parameters and the type the table records.
\end{lemma}

\begin{proof} \leanok
\uses{def:Witness-ReifiedDecl-Pinned, def:Witness-SourceTableAdequate, def:TableSafe}
A successful lookup exhibits a member of the declaration column, so the \code{decls} field of
adequacy applies and its pinning half gives the environment constant with the two equalities; the
safety clause of \code{TableSafe} is read at the same lookup. This is the single entry point from a
decidable table lookup to environment facts.
\end{proof}

\begin{lemma}[The rung's subject translates]\label{lem:Witness-trExprS_const_of_table}
\lean{LeanToLambdaBox.Witness.trExprS_const_of_table} \leanok
\uses{def:Witness-SourceTableAdequate, def:Witness-SourceTable, def:TableSafe, def:ErasureSpec, def:lean4lean-grounding}
Under the specification bundle, an adequate and safe table, and the computed side condition that the
table records no level parameters for a name, the source constant of that name translates to the
model constant of the same name, at any modelled local context.
\end{lemma}

\begin{proof} \leanok
\uses{lem:Witness-tableDeclPin, lem:Witness-trExprS_const}
Case on the table lookup; the \code{none} branch contradicts the side condition. In the \code{some}
branch the pinning lemma supplies the environment constant, and the side condition transports the
table's level-parameter list to the constant's, so the monomorphic instance of the translation lemma
applies. This is what discharges the capstone's \code{hwt} at all eight rungs, its one side
condition being \code{by rfl} on the reified table; the same route is exercised outside the ladder
at \code{Nat.add}.
\end{proof}

\begin{lemma}[The tabled subject's type, in the model]\label{lem:Witness-hasType_const_of_table}
\lean{LeanToLambdaBox.Witness.hasType_const_of_table} \leanok
\uses{def:Witness-SourceTableAdequate, def:Witness-SourceTable, def:TableSafe, def:ErasureSpec, def:lean4lean-grounding}
With the same hypotheses, and the further computed condition that the table records the subject's
type as a constant, the model types the subject at that constant.
\end{lemma}

\begin{proof} \leanok
\uses{lem:Witness-tableDeclPin, lem:Witness-hasType_const_nil}
As in the previous lemma, with the type equation transported alongside the level-parameter equation
and the monomorphic typing lemma applied in place of the translation lemma. Rungs whose subject is
function-typed fall outside this lemma; see Lemma~\ref{lem:Green-g8_trExprS_spine}.
\end{proof}

\begin{lemma}[The checker route: a successful run reports a translation]\label{lem:Witness-trTyping_of_checkType_run}
\lean{LeanToLambdaBox.Witness.trTyping_of_checkType_run, LeanToLambdaBox.Witness.trTyping_of_checkType_run_closed} \leanok
\uses{def:ErasureSpec, def:lean4lean-grounding}
If lean4lean's executable type checker, run purely at a well-formed environment and a well-formed
ambient local context whose free variables the kernel's fixed name generator reserves, returns a
type for a term whose free variables lie in that context, then both the term and the reported type
translate at that context and the model types the first at the second. In the closed instance, taken
at the specification bundle's own model connection, a closed source term the checker accepts at the
empty local context translates together with its reported type in the model environment.
\end{lemma}

\begin{proof} \leanok
\uses{def:lean4lean-grounding, lem:Lean4Lean-TypeChecker-M-WF-run-prime, def:Lean4Lean-TypeChecker-VContext-ofMLCtx, asm:ErasureSpec-env_connect, asm:lean4lean-trust, asm:lean-reflection-axioms}
lean4lean's own well-formedness lemma for \code{checkType} states the translation and its typing at
an abstract verified context; this repository's checker adequacy supplies the initial state at an
ambient multi-level local context, so a pure run returning \code{ok} witnesses both translations and
the typing. The closed instance instantiates the ambient context with the empty one, obtaining the
model connection from the bundle's \code{env\_connect} field.
\end{proof}

\begin{remark}
This is the second route to the capstone's translation binders, for subjects that are not constants:
it replaces a translation binder by an executable run equation of the same kind as \code{hrun}.
Unlike the table route it goes through lean4lean's executable checker layer, and its measured
footprint accordingly contains \code{sorryAx} and the reflection axioms. \textbf{No rung uses it}
--- the ladder's eight subjects are all constants.
\end{remark}

\section{The capstone theorem}
\label{sec:capstone:theorem}

\begin{theorem}[The shipping erasure is correct at first-order answers]\label{thm:shipping_erase_correct_firstorder}
\lean{LeanToLambdaBox.shipping_erase_correct_firstorder} \leanok
\uses{def:Witness-SourceTableAdequate, def:TableSafe, def:TableBlocks, def:ConfigPinned, def:CompilerBodies, def:Supported, def:Erases, def:ErasesEnv, def:Lower, def:LowerEnv, def:SEval, def:SEvalFlags, def:WcbvEval, def:eraseFlags, def:NoBox, def:FirstOrderInd, def:LBWfPeregrine, def:NoBodylessRefs, def:SpecEnv, def:ErasureBridge, def:Erasure-erase, def:Erasure-prepare_erasure, def:Erasure-visitExpr, def:ErasureSpec, def:EraserAsks, def:UpstreamAsks, def:lean4lean-grounding, asm:htbl, asm:hsafe, asm:hblk, asm:hcb, asm:hprep, asm:hrun, asm:ErasureBridge-erasesEnv, asm:ErasureBridge-lowerEnv}
Fix a Lean environment \code{lenv}, a model environment \code{env} and a name-generator reading, a
source table, an erasure configuration \code{cfg}, source terms $e$ and $pe$, an emitted environment
$\Gamma$ and an emitted term $t$. Assume
\begin{enumerate}
\item $P$: the specification bundle \code{ErasureSpec} for Lean's elaboration primitives and the
    model connection (eight fields);
\item $E$: this repository's own obligations \code{EraserAsks} (five fields);
\item $A$: the four load-bearing upstream obligations \code{UpstreamAsks};
\item \code{htbl}: the table is adequate for \code{lenv};
\item \code{hsafe}: every tabled declaration is safe in \code{lenv} (\code{TableSafe});
\item \code{hblk}: every block the run installs a fix-variable map for is tabled, lambda-headed and
    non-erasable (\code{TableBlocks});
\item \code{hcfg}: the configuration is the pinned one (\code{ConfigPinned});
\item \code{hcb}: every tabled compiler body is kernel-typeable at its declared type
    (\code{CompilerBodies});
\item \code{hwt}: the prepared term $pe$ translates into the model;
\item \code{hsup}: $pe$ lies in the supported fragment for the table;
\item \code{hprep}: preparation run on $e$ from the empty state at \code{cfg} returns $pe$ and the
    empty state;
\item \code{hrun}: \code{Erasure.erase e cfg} returns the untyped program $(\Gamma, t)$ together
    with an inlining list, which the theorem binds and discards;
\item \code{hnb}: every constant $(\Gamma,t)$ reaches is declared with a body
    (\code{NoBodylessRefs});
\item \code{hwf}: $(\Gamma,t)$ satisfies \code{LBWfPeregrine};
\item \code{hbridge}: for the state the \code{visitExpr} run from $pe$ ends in there is a
    specification environment $\Gamma_{spec}$ with \code{SpecEnv}, such that for every $t_0$ with
    $\Erases\ pe\ t_0$ and $\Lower\ \Gamma_{spec}\ t_0\ t$, the pair $(\Gamma_{spec}, \Gamma)$
    satisfies \code{ErasureBridge} at $t_0$.
\end{enumerate}
Then there are a specification environment $\Gamma_{spec}$ and a \lbox{} term $t_0$ such that: (a)
the preparation ran as assumed; (b) $pe$ erases to $t_0$; (c) $\Gamma_{spec}$ erases the source
environment along everything $t_0$ reaches; (d) the lowering takes $t_0$ to $t$; (e) it takes
$\Gamma_{spec}$ to $\Gamma$; (f) $(\Gamma,t)$ satisfies \code{LBWfPeregrine}; and (g) the
\textbf{observable}: for every list of source arguments and every list of \lbox{} arguments of the
same length, each of which carries an erasure of the corresponding source argument, its lowering and
its own environment clause, if the applied spine translates, if the source evaluation of $e$ applied
to the arguments yields a value $v$, if $v$ translates and the model types its translation at a
spine of an inductive $I$ that is first order in the model, then $v$ has a \textbf{unique} erasure
$tv_0$, its lowering $tv$ is \textbf{box-free}, and the emitted program evaluates $t$ applied to the
\lbox{} arguments to that very $tv$ under \lbox{}'s call-by-value semantics.

Beyond the prose, the Lean statement fixes both evaluations at flag sets --- the source side at
\code{fullFlags}, the target side at \code{eraseFlags} --- keeps the local contexts and level scopes
empty throughout, and states the argument-indexed premise with \code{getElem!} on both lists.
\end{theorem}

\begin{proof} \leanok
\uses{thm:erase_run_ok, thm:erasure_bridge_of_run, thm:prepare_sound, thm:erases_correct, thm:firstorder_erases_core, lem:noBox_lower_of_foSpine, lem:Erases-mkApps, lem:Lower-target_fix, lem:ErasesEnv-mkApps, asm:lean4lean-trust, asm:lean-reflection-axioms}
The run-split theorem factors the shipping run into a preparation and a \code{visitExpr} run from
the empty state and identifies the emitted term with the one the statement names; injecting
\code{hprep} against that split fixes $pe$, the empty state and the world token. Applying
\code{hbridge} to the \code{visitExpr} run supplies the specification environment and, once the
bridge theorem has produced the intermediate term $t_0$ with its erasure and its lowering, the two
environment results. For the observable, a generic list lemma turns the per-argument premise into a
list of erasures of the arguments; the preparation-soundness theorem carries the source evaluation
from $e$ to $pe$ under the spine; the three spine lemmas for $\Erases$, $\Lower$ and $\ErasesEnv$
assemble the corresponding spine facts; \code{erases\_correct} is applied \textbf{once}, at the
spine, and yields the target evaluation; and the first-order core gives the answer's
constructor-spine shape and the uniqueness of its erasure, from which the box-freedom lemma turns
that shape into box-freedom of the lowered value.
\end{proof}

The fifteen binders fall into classes. $P$, $E$ and $A$ are hypothesis bundles whose fields
Chapter~\ref{chap:trust} lists one by one. \code{htbl}, \code{hsafe}, \code{hblk}, \code{hcb},
\code{hprep} and \code{hrun} are statements about the live environment or about a monadic run; some
are mechanised outside the kernel, none is a theorem. \code{hcfg}, \code{hsup}, \code{hnb} and
\code{hwf} are decidable at a concrete program and are discharged at every rung. \code{hwt} is
discharged by the table route. \code{hbridge} is the hole.

\begin{remark}
Two conjuncts of the conclusion are literally hypotheses of the theorem: conjunct (a) is the binder
\code{hprep} verbatim, and conjunct (f) is the binder \code{hwf} verbatim. Both are carried, not
derived. The binder \code{hnb} is the theorem's one genuinely unused binder: it appears nowhere in
the proof term and is not among the conclusion's conjuncts. The binder \code{hblk} is forwarded to
the bridge theorem, but \code{TableBlocks} is consumed by no theorem in the closure: its only
consumer in the tree is a lemma that is itself consumed by nothing.
\end{remark}

\begin{remark}
Explicitly \textbf{not} concluded: the expanded-fixpoint condition corresponding to MetaRocq's
\code{expanded\_eprogram}. The eraser emits bare fixpoint constant bodies, so the stronger
precondition is false of its output --- shipping finding F-ETA. The theorem concludes the weaker
\code{LBWfPeregrine}, which the output does satisfy. Also: the proof derives the identification of
the emitted environment with the run's final state and then discards it, keeping only the term
component, so the statement relates $\Gamma$ to the run only through \code{hrun}, and
\code{hbridge}'s lowering field speaks about the \emph{given} $\Gamma$.
\end{remark}

\begin{remark}
Correspondence with [S, Sec. 7.3--7.4]: the same shape with three deliberate substitutions ---
\lbox{}'s own weak call-by-value evaluation as target semantics, the extra relation $\Lower$ (and
hence the split of the environment side into $\ErasesEnv$ and $\LowerEnv$), and a source evaluation
reading the compiler-body table rather than the kernel's $\delta$ rules (tag N8). The first-order
restriction is where MetaRocq's own first-order inductive machinery sits, and \code{NoBodylessRefs}
is this repository's counterpart of \code{axiom\_free}. Finally, \code{hbridge} is inhabited
\textbf{nowhere}: \code{bridgeEnv\_of\_regInv} (Chapter~\ref{chap:coldstart}) is the proved
interface that would derive its two fields from a registration invariant, but no theorem establishes
that invariant for a run of the shipping eraser.
\end{remark}

\section{The ladder's data}
\label{sec:capstone:data}

\begin{definition}[The eight rung subjects]\label{def:spikeZero}
\lean{spikeZero, spikeLit, spikeLet, spikeProj, spikeCase, spikeRec, spikeFix, arithClosed, benchArith} \leanok
The ladder's subjects, in order: \code{spikeZero} is \code{Nat.zero}; \code{spikeLit} is
\code{Nat.succ 3}; \code{spikeLet} is \code{let x := 2; Nat.succ x}; \code{spikeProj} is
\code{(Prod.mk 1 2).1}; \code{spikeCase} is a \code{Nat.casesOn} applied to the constructor value
\code{Nat.succ (Nat.succ Nat.zero)}, whose nullary branch is a \emph{thunk application} --- the
shape Lean's \code{match} compiler gives a nullary alternative --- and which is spelled with
\code{Nat.casesOn} rather than \code{match} so that the iota redex sits in the subject rather than
behind a matcher constant; \code{spikeRec} is the analogous recursive definition and \code{spikeFix}
its application \code{spikeRec (Nat.succ (Nat.succ Nat.zero))}; \code{arithClosed} is
\code{benchArith 0}; and \code{benchArith n} is the tracked benchmark program $2^{((n \cdot 3) - n)
+ 3}$ itself.
\end{definition}

\begin{remark}
All of \code{spikeZero} through \code{arithClosed} are declared at the \textbf{root} namespace,
deliberately: their kernames are then \code{rootKername "spikeZero"} and so on, which is what the
committed emitted environments record. \code{benchArith} lives in the frozen benchmark source
\code{VerifyBench/Src/Arith.lean}, held byte-identical against its expected copy and against the
sibling \code{benchmarks} repository by \code{scripts/frozen.sh}.
\end{remark}

\begin{definition}[The pinned configuration]\label{def:Green-spikeConfig}
\lean{LeanToLambdaBox.Green.spikeConfig} \leanok
\uses{def:Erasure-ErasureConfig, def:ConfigPinned}
Every rung is erased under one configuration: \code{extern := .preferLogical}, \code{nat := .peano},
\code{csimp := false}, the remaining fields at their defaults.
\end{definition}

\begin{remark}
This is \textbf{not} the default \code{\#erase} configuration: \code{ConfigPinned} differs from the
default on three of its five fields (\code{extern}, \code{nat} and \code{csimp}). Scope rows N1--N4
are pinned here, in particular \code{csimp} off and constructor-argument pruning off.
\end{remark}

\begin{lemma}[The rungs' configuration is the pinned one]\label{lem:Green-spike_configPinned}
\lean{LeanToLambdaBox.Green.spike_configPinned} \leanok
\uses{def:Green-spikeConfig, def:ConfigPinned}
The rung configuration satisfies \code{ConfigPinned}.
\end{lemma}

\begin{proof} \leanok
\uses{def:Green-spikeConfig, def:ConfigPinned}
\code{ConfigPinned} is a conjunction of five field equations, each true by \code{rfl} of the literal
configuration. This discharges the capstone's \code{hcfg} at all eight rungs.
\end{proof}

\begin{definition}[The literal answers: peano numerals in \lbox{}]\label{def:Green-peanoLB}
\lean{LeanToLambdaBox.Green.peanoLB, LeanToLambdaBox.Green.litLB, LeanToLambdaBox.Green.natIid, LeanToLambdaBox.Green.ofNatIid, LeanToLambdaBox.Green.ofNatKn, LeanToLambdaBox.Green.instOfNatNatKn} \leanok
\uses{def:LBTerm, def:InductiveId, def:Kername}
\code{peanoLB n} is the \lbox{} peano numeral: the zero constructor node under $n$ applications of
the successor constructor node, each constructor node carrying an \emph{empty} argument list. Under
\code{nat := .peano} this is what the eraser emits for a \code{Nat} literal, so it is the language
every rung's committed answer is written in. \code{litLB n} is the emitted image of a \emph{source}
\code{Nat} literal: the class projection \code{OfNat.ofNat} applied to the erased carrier type --- a
$\Box$ --- then to the numeral and then to the instance, because the source syntax \code{3} is
\code{@OfNat.ofNat Nat 3 (instOfNatNat 3)}.
\end{definition}

\begin{remark}
Two things are visible in this data. First, the \emph{applied} constructor convention: every emitted
constructor node carries an empty argument list (measured 982 of 982 across the five benchmark
programs, finding F-ETA2), as opposed to CertiCoq's block form. Second, the box rule of [L] and [S,
Sec. 7.3] in the committed output: \code{litLB} shows a $\Box$ where the carrier type was.
\end{remark}

\begin{definition}[The per-rung committed data]\label{def:Green-g1Table}
\lean{LeanToLambdaBox.Green.g1Table, LeanToLambdaBox.Green.eG1, LeanToLambdaBox.Green.g1Env, LeanToLambdaBox.Green.g1Term, LeanToLambdaBox.Green.g1Answer, LeanToLambdaBox.Green.g2Table, LeanToLambdaBox.Green.g3Table, LeanToLambdaBox.Green.g4Table, LeanToLambdaBox.Green.g5Table, LeanToLambdaBox.Green.g6Table, LeanToLambdaBox.Green.g7Table, LeanToLambdaBox.Green.g8Table, LeanToLambdaBox.Green.g8Arg, LeanToLambdaBox.Green.arithTower, LeanToLambdaBox.Green.g7Tag, LeanToLambdaBox.Green.g8Tag} \leanok
\uses{def:Witness-elabReify, def:Green-peanoLB, def:GlobalDeclarations, def:LBTerm}
Each rung commits four literals and one answer. The table \code{g<i>Table} is \code{reify\%
<subject>}, the reified slice of the environment this module elaborates in. The subject \code{eG<i>}
is the source term \code{\#erase <subject>} elaborates, namely the constant of the subject's name.
The environment \code{g<i>Env} and the term \code{g<i>Term} are the emitted program, transcribed by
hand from the committed \code{.ast} file under \code{VerifyBench/ast/Spikes/}. The answer
\code{g<i>Answer} is the literal peano numeral the rung claims. G7 and G8 share \code{arithTower},
thirty-nine declarations --- \code{benchArith}'s body, the power, multiplication, subtraction and
addition classes with their \code{Nat} and homogeneous-bridge instances, the four \code{Nat}
operations as fixpoint blocks plus \code{Nat.pred}, and \code{PUnit}, \code{Nat}, \code{OfNat} and
\code{instOfNatNat}. \code{g8Env} \textbf{is} \code{arithTower g8Tag}, since \code{benchArith} is
itself the tower's first entry; \code{g7Env} conses one further declaration, \code{arithClosed}'s
own body, onto the same tower --- parametrised by the macro-scope tag Lean's matcher inliner writes
into the ten emitted \code{let} binders, which is why the two rungs are transcribed at two
different tags.
\end{definition}

\begin{remark}
This is \textbf{data}, not a claim. Its correspondence to a real run is the binder \code{hrun},
byte-diffed by \code{lake exe green-check}; no Lean proof can exist, \code{IO.RealWorld} being
opaque. The transcription includes hygienic binder names and macro-scope tags, so it is a
maintenance risk guarded only by that external diff. One visible asymmetry is worth recording: the
emitted \code{Nat.add} and \code{Nat.mul} fixpoints recurse on their \textbf{second} argument, the
Lean-versus-Rocq difference the benchmark repository also documents.
\end{remark}

\section{The hypotheses a computation settles}
\label{sec:capstone:checked}

The four lemmas of this section exist once per rung; each is stated below for G1, and the seven
analogues differ only in the committed literals they are read at.

\begin{lemma}[Each subject is inside the fragment]\label{lem:Green-g1_supported}
\lean{LeanToLambdaBox.Green.g1_supported, LeanToLambdaBox.Green.g2_supported, LeanToLambdaBox.Green.g3_supported, LeanToLambdaBox.Green.g4_supported, LeanToLambdaBox.Green.g5_supported, LeanToLambdaBox.Green.g6_supported, LeanToLambdaBox.Green.g7_supported, LeanToLambdaBox.Green.g8_supported} \leanok
\uses{def:supportedB, def:Green-g1Table}
For each rung, the Boolean fragment checker returns success on the rung's reified table at depth
fuel 8 and the rung's subject.
\end{lemma}

\begin{proof} \leanok
\uses{def:supportedB, def:Green-g1Table}
The checker's success is decided by the \textbf{kernel} (\code{by decide +kernel}) on the
\code{isOk} projection, then a case analysis on the \code{Except} value turns the Boolean verdict
into the equation. The kernel is needed because \code{String.isPrefixOf}, which the matcher-name
test calls, does not reduce in the elaborator. Combined with the checker's soundness theorem this
discharges the capstone's \code{hsup} at every rung.
\end{proof}

\begin{lemma}[Every reached constant has a body]\label{lem:Green-g1_noBodylessRefs}
\lean{LeanToLambdaBox.Green.g1_noBodylessRefs, LeanToLambdaBox.Green.g2_noBodylessRefs, LeanToLambdaBox.Green.g3_noBodylessRefs, LeanToLambdaBox.Green.g4_noBodylessRefs, LeanToLambdaBox.Green.g5_noBodylessRefs, LeanToLambdaBox.Green.g6_noBodylessRefs, LeanToLambdaBox.Green.g7_noBodylessRefs, LeanToLambdaBox.Green.g8_noBodylessRefs} \leanok
\uses{def:NoBodylessRefs, def:Green-g1Table}
For each rung, every constant the committed emitted program reaches is declared with a body.
\end{lemma}

\begin{proof} \leanok
\uses{def:NoBodylessRefs, def:Green-g1Table}
\code{decide +kernel} on the committed literals; at G7 this is forty declarations and a forty-round
reachability closure, the ladder's longest kernel computation.
\end{proof}

\begin{remark}
This discharges \code{hnb}. Since \code{hnb} is an unspent binder of the capstone, the honest
statement is that the rungs settle a hypothesis the theorem does not use: it is the exclusion
mechanism for the body-less \code{Eq.rec} finding F-EQREC at the \emph{program} level, not a premise
of the theorem's proof. Among the five benchmark programs only Fannkuch fails it.
\end{remark}

\begin{lemma}[Each emitted program is well formed for peregrine]\label{lem:Green-g1_wf}
\lean{LeanToLambdaBox.Green.g1_wf, LeanToLambdaBox.Green.g2_wf, LeanToLambdaBox.Green.g3_wf, LeanToLambdaBox.Green.g4_wf, LeanToLambdaBox.Green.g5_wf, LeanToLambdaBox.Green.g6_wf, LeanToLambdaBox.Green.g7_wf, LeanToLambdaBox.Green.g8_wf} \leanok
\uses{def:LBWfPeregrine, def:Green-g1Table}
For each rung, the committed emitted program satisfies \code{LBWfPeregrine}, the twelve clauses
peregrine's first pass reads of a \lbox{} program.
\end{lemma}

\begin{proof} \leanok
\uses{thm:lbWfPeregrine_of_check, def:LBWfPeregrine, def:Green-g1Table}
The Boolean checker decides all twelve clauses on the committed environment and term in the kernel,
and the transport theorem turns its verdict into the structure. The eight kernel checks cost about
twelve seconds of elaboration, most of it at G7 and G8.
\end{proof}

\begin{remark}
A recently corrected vacuity belongs here. Until the closing round of work the binder-name clause of
\code{LBWfPeregrine} demanded the alphanumeric identifier class that the kername constructor applies
to \emph{kernames}, which emitted \emph{binder} names do not satisfy (measured 0/1/1/1/0/0/34/34
offending names at G1 through G8), so G2, G3, G4, G7 and G8 had an unsatisfiable hypothesis and were
vacuous. The clause now reads \code{PrintableBinders}, under which the count is 0 at all eight
rungs. Any older description of the ladder is stale on this point.
\end{remark}

\begin{lemma}[Each emitted program evaluates to its literal answer]\label{lem:Green-g1_eval}
\lean{LeanToLambdaBox.Green.g1_eval, LeanToLambdaBox.Green.g2_eval, LeanToLambdaBox.Green.g3_eval, LeanToLambdaBox.Green.g4_eval, LeanToLambdaBox.Green.g5_eval, LeanToLambdaBox.Green.g6_eval, LeanToLambdaBox.Green.g7_eval, LeanToLambdaBox.Green.g8_eval} \leanok
\uses{def:WcbvEval, def:eraseFlags, def:Green-g1Table, def:Green-peanoLB}
For each rung the committed emitted program evaluates, under \lbox{}'s call-by-value semantics at
the erasure flag set, to the committed literal peano answer. At G8 the evaluated term is the emitted
term \textbf{applied to} the \lbox{} numeral 0.
\end{lemma}

\begin{proof} \leanok
\uses{thm:lbEval_sound, def:WcbvEval, def:Green-peanoLB}
The certified evaluator is run in the kernel at a fixed fuel --- 8, 16, 16, 16, 16, 24, 64, 64 at G1
through G8 --- and its soundness theorem transports the computed value to a derivation of the
evaluation relation. This is the ladder's anti-vacuity device: together with determinism of \lbox{}
evaluation it forces the capstone's existential target value to be the literal numeral, so the
conclusion cannot be met by $\Box$ or by a stuck term. The evaluator's own measured footprint is
\code{[propext]}, the cleanest in the ledger.
\end{proof}

\begin{lemma}[G1's compiler bodies are typed]\label{lem:Green-g1_compilerBodies}
\lean{LeanToLambdaBox.Green.g1_compilerBodies, LeanToLambdaBox.Green.g1Decls, LeanToLambdaBox.Green.g1_decls_eq, LeanToLambdaBox.Green.g1_body?_some, LeanToLambdaBox.Green.g1_pinned_subject, LeanToLambdaBox.Green.g1_pinned_zero} \leanok
\uses{def:CompilerBodies, def:Witness-SourceTableAdequate, def:TableSafe, def:ErasureSpec, def:Green-g1Table}
Under the specification bundle and an adequate, safe G1 table, \code{CompilerBodies} holds of G1's
body column: the table's declaration column is the subject together with the two \code{Nat}
constructors, only the subject carries a body, and that body is the constant \code{Nat.zero}, which
the pinned table types at the subject's declared type in the model.
\end{lemma}

\begin{proof} \leanok
\uses{def:Witness-SourceTableAdequate, asm:ErasureSpec-decl_adequate, def:CompilerBodies}
The declaration column is settled by computation --- the equation identifying it with the literal
list is \code{rfl} --- so the only tabled body is the subject's, and it is the constant
\code{Nat.zero}. Adequacy pins the subject and \code{Nat.zero} in the ambient environment,
monomorphically and at type \code{Nat}. The bundle's \code{decl\_adequate} field then supplies the
model constant for \code{Nat.zero} and the translation of its declared type; inverting that
translation gives the model typing of the body at the subject's type, and a constant-unfolding step
supplies the required definitional equality.
\end{proof}

\begin{remark}
This is \textbf{the only rung at which \code{hcb} is discharged.} It stays a binder at G2 through
G8: the tables of G2--G4 carry the class projection \code{OfNat.ofNat}, whose translation routes
through lean4lean's projection layer, which is entirely unproven at the pin; G5 and G6 carry an
eliminator spine and a recursive body; and G7/G8's Arith table has ten projection-bearing bodies
among thirty, with twenty-eight of its forty-four declarations polymorphic.
\end{remark}

\section{G5: the source evaluation, derived}
\label{sec:capstone:g5}

\begin{definition}[G5's source-side vocabulary]\label{def:Green-peanoSrc}
\lean{LeanToLambdaBox.Green.peanoSrc, LeanToLambdaBox.Green.g5Motive, LeanToLambdaBox.Green.g5Minor0, LeanToLambdaBox.Green.g5Minor1, LeanToLambdaBox.Green.g5Body} \leanok
\uses{def:Green-g1Table}
A peano numeral in \emph{source constructor form}; G5's motive; its thunked nullary minor premise,
the identity thunk applied to \code{Unit.unit}; its successor minor premise, the identity function;
and \code{g5Body}, the eliminator spine \code{Nat.casesOn} applied to the motive, the constructor
value 2 and the two minors --- which is the body G5's table records for the subject.
\end{definition}

\begin{remark}
The thunked nullary branch is not an artefact of the fixture: it is the shape Lean's \code{match}
compiler gives a nullary alternative, and it is exactly what makes the iota rule's per-branch
obligation (restriction N20) visible.
\end{remark}

\begin{lemma}[G5's tabled bodies, computed]\label{lem:Green-g5_body_eq}
\lean{LeanToLambdaBox.Green.g5_body_eq, LeanToLambdaBox.Green.g5_unit_body_eq} \leanok
\uses{def:Green-peanoSrc, def:Green-g1Table}
G5's table records the eliminator spine as the subject's body and \code{PUnit.unit} as
\code{Unit.unit}'s.
\end{lemma}

\begin{proof} \leanok
\uses{def:Green-peanoSrc, def:Green-g1Table}
Both are \code{rfl} on the reified literal.
\end{proof}

\begin{definition}[What the model must declare about \code{Nat}]\label{def:Green-SpikeNatFacts}
\lean{LeanToLambdaBox.Green.SpikeNatFacts, LeanToLambdaBox.Green.SpikeNatFacts.natInd, LeanToLambdaBox.Green.SpikeNatFacts.natZero, LeanToLambdaBox.Green.SpikeNatFacts.natSucc, LeanToLambdaBox.Green.SpikeNatFacts.natCases, LeanToLambdaBox.Green.SpikeNatFacts.natCasesOrigin, LeanToLambdaBox.Green.SpikeNatFacts.natInf} \leanok
\uses{def:IndInfo, def:CtorOf, def:CasesOnShape, def:ConstOrigin, def:InformativeInd}
A bundle of six model-side facts: \code{Nat} is declared with no parameters and two constructors of
zero and one field; \code{Nat.zero} and \code{Nat.succ} are its constructors at indices 0 and 1;
\code{Nat.casesOn} has one argument before the major premise and two minors --- the segmentation at
which the iota rule splits a spine; \code{Nat.casesOn} is declared as a \emph{plain} constant; and
\code{Nat} eliminates into data, its result sort never being \code{Prop}.
\end{definition}

\begin{remark}
Each field is a datum the reified table records for the ambient environment, \emph{read in the
model}; the transport between the two is upstream ask 3, which is why this is a bundle of hypotheses
rather than a theorem. The informativeness field is restriction N18.
\end{remark}

\begin{lemma}[The \code{Nat} bundle is satisfiable]\label{lem:Green-spikeNatFacts_natEnv}
\lean{LeanToLambdaBox.Green.spikeNatFacts_natEnv} \leanok
\uses{def:Green-SpikeNatFacts, prop:NatWitness-seval_iota_fires}
The \code{Nat} bundle holds at the model fixture of the source-evaluation chapter and its inductive
identifier.
\end{lemma}

\begin{proof} \leanok
\uses{def:Green-SpikeNatFacts, prop:NatWitness-seval_iota_fires}
The source-evaluation fixture builds a well-formed model environment declaring \code{Nat}, its
constructors, its recursor and \code{Nat.casesOn}; every field of the bundle is one of that
fixture's own theorems, so the bundle is assembled directly. The measured footprint is
\code{[propext, Classical.choice, Quot.sound]}: no \code{sorryAx}. This is the non-vacuity guard for
G5's and G8's model-side bundle.
\end{proof}

\begin{definition}[What the model must declare about \code{PUnit}]\label{def:Green-SpikeUnitFacts}
\lean{LeanToLambdaBox.Green.SpikeUnitFacts, LeanToLambdaBox.Green.SpikeUnitFacts.punitInd, LeanToLambdaBox.Green.SpikeUnitFacts.punitUnit} \leanok
\uses{def:IndInfo, def:CtorOf}
\code{PUnit} is declared with no parameters and one nullary constructor, and \code{PUnit.unit} is
that constructor.
\end{definition}

\begin{remark}
This bundle is kept separate from the \code{Nat} one because the model fixture declares a single
block, so the two halves would be satisfiable only at different fixtures. \textbf{This half is
inhabited nowhere}: G5's claim that its source evaluation is not assumed therefore holds modulo a
bundle half whose consistency is not exhibited.
\end{remark}

\begin{lemma}[Source peano numerals are values]\label{lem:Green-peanoSrc_eval}
\lean{LeanToLambdaBox.Green.peanoSrc_eval} \leanok
\uses{def:Green-SpikeNatFacts, def:Green-peanoSrc, def:SEval}
Given the \code{Nat} bundle, every peano numeral in source constructor form evaluates to itself, at
any body column, level scope and flag set.
\end{lemma}

\begin{proof} \leanok
\uses{def:Green-SpikeNatFacts, def:Green-peanoSrc, def:SEval}
Induction on the numeral. At zero and at a successor the spine is within the constructor's arity, so
the constructor-value rule of the source evaluation applies, with the inductive hypothesis
discharging the single argument in the successor case.
\end{proof}

\begin{lemma}[Neither the subject nor the thunk's argument is an eliminator]\label{lem:Green-not_casesOn_spikeCase}
\lean{LeanToLambdaBox.Green.not_casesOn_spikeCase, LeanToLambdaBox.Green.not_casesOn_unitUnit} \leanok
\uses{def:CasesOnShape}
No \code{CasesOnShape} holds of \code{spikeCase}, nor of \code{Unit.unit}, at any inductive type,
prefix depth or minor count.
\end{lemma}

\begin{proof} \leanok
\uses{def:CasesOnShape}
The first component of \code{CasesOnShape} is a decidable condition on the name, refuted by
\code{decide}. These are the side conditions the delta rule of the source evaluation takes; they are
what makes a delta step the only available move at those two heads.
\end{proof}

\begin{proposition}[G5's source evaluation, derived]\label{prop:Green-g5_seval}
\lean{LeanToLambdaBox.Green.g5_seval} \leanok
\uses{def:Green-SpikeNatFacts, def:Green-SpikeUnitFacts, def:Green-peanoSrc, def:Green-g1Table, def:SEval, def:StepDefeq, def:SEvalFlags}
Given the \code{Nat} and \code{PUnit} bundles and three definitional-step facts --- that the subject
is definitionally its tabled body, that \code{Unit.unit} is definitionally \code{PUnit.unit}, and
that the eliminator spine is definitionally the selected minor applied to the constructor's field
--- the source evaluation takes G5's subject to the source numeral 1, reading the body column off
G5's reified table at the full flag set.
\end{proposition}

\begin{proof} \leanok
\uses{lem:Green-g5_body_eq, lem:Green-peanoSrc_eval, lem:Green-not_casesOn_spikeCase, def:Green-SpikeNatFacts, def:Green-SpikeUnitFacts, def:SEval}
A delta step at the subject --- licensed by the tabled body, by the subject not being an eliminator,
and by the first definitional-step fact --- exposes the eliminator spine. The iota rule then fires
at \code{Nat.casesOn}: the discriminant is a constructor value by the numeral lemma, the block data
comes from the \code{Nat} bundle, and the selected minor applied to the constructor's field
beta-reduces to the numeral 1. The rule also owes a derivation for the \emph{unselected} nullary
minor, which is a thunk application; that is discharged by a beta step whose argument needs its own
delta step at \code{Unit.unit}, using the \code{PUnit} bundle and the second definitional-step fact.
\end{proof}

\begin{remark}
This is the only place in the repository where the capstone's source-evaluation premise is a term
rather than a binder, hence the ladder's non-vacuity high-water mark on the source side. Its
measured footprint is \code{[propext, Classical.choice, Quot.sound]}: no \code{sorryAx}. Design tags
N20 (the per-branch obligation of the iota rule) and N18 (informativeness). No source-evaluation
witness is constructed at Arith: \code{benchArith 0} runs 45 recursive calls, each a delta on a
fixpoint-carrying constant and an iota on its compiled match, and N20 would owe a derivation for
every unselected minor of each --- a recorded fallback, not an oversight.
\end{remark}

\section{G8: the applied rung}
\label{sec:capstone:g8}

\begin{lemma}[G8's spine translates]\label{lem:Green-g8_trExprS_spine}
\lean{LeanToLambdaBox.Green.g8_trExprS_spine} \leanok
\uses{def:Witness-SourceTableAdequate, def:Green-g1Table, def:TableSafe, def:ErasureSpec, def:lean4lean-grounding}
Under the specification bundle and an adequate, safe G8 table, the source application of
\code{benchArith} to \code{Nat.zero} translates to the model application of the two model constants.
\end{lemma}

\begin{proof} \leanok
\uses{lem:Witness-trExprS_const_of_table, lem:Witness-hasType_const_of_table, lem:Witness-tableDeclPin, asm:ErasureSpec-decl_adequate}
The application arm of the translation relation needs the function's arrow type in the model and the
argument's typing. The table records \code{benchArith}'s type as a dependent arrow between two
\code{Nat} constants and records no level parameters, both by \code{rfl}; the pinning lemma
transports this to the ambient environment, and inverting \code{decl\_adequate}'s translation of
that arrow twice --- once at the domain, once at the codomain, using that a translated constant is
the model constant of the same name --- identifies the model type. The two components then come from
the table route and the argument's typing from the tabled-type lemma.
\end{proof}

\begin{remark}
The capstone's observable reads its translation premise at the \emph{applied spine}, so a
function-typed subject cannot use the table route alone. G8 is the only rung with a non-empty spine.
\end{remark}

\begin{lemma}[The argument's image reaches one key]\label{lem:Green-reachableFrom_peanoLB_zero}
\lean{LeanToLambdaBox.Green.reachableFrom_peanoLB_zero} \leanok
\uses{def:Green-peanoLB, def:ReachableFrom, def:LBTerm-envLookup}
If a specification environment answers \code{Nat}'s block key with an inductive declaration, then
the only kername reachable from the \lbox{} numeral 0 in that environment is that key.
\end{lemma}

\begin{proof} \leanok
\uses{def:Green-peanoLB, def:ReachableFrom, def:LBTerm-envLookup}
The numeral 0 is a nullary constructor node whose constant references are the single block key. The
reachability closure is computed by iterating an expansion step; an induction on the iteration count
shows the singleton is a fixed point, because a block key answers with an inductive declaration,
which the expansion does not unfold. Membership in the singleton is then the claimed equation.
\end{proof}

\begin{lemma}[G8's argument is covered by the subject's environment]\label{lem:Green-g8_argErasesEnv}
\lean{LeanToLambdaBox.Green.g8_argErasesEnv} \leanok
\uses{def:Green-SpikeNatFacts, def:Green-peanoLB, def:Green-g1Table, def:ErasesEnv, def:ReachableFrom}
Given the \code{Nat} bundle, an environment-erasure of the subject's erasure, and the fact that the
subject's erasure reaches \code{Nat}'s block key in that environment, the same environment erases
the source environment along the argument's \lbox{} image.
\end{lemma}

\begin{proof} \leanok
\uses{lem:Green-reachableFrom_peanoLB_zero, def:Green-SpikeNatFacts, def:ErasesEnv}
The \code{Nat} block's declaration in the specification environment is obtained from the subject's
environment-erasure at the reachability fact. Of the seven conditions of $\ErasesEnv$, two --- the
key clause and the tabled clause --- are reused verbatim from the subject's derivation, and the
remaining five are re-derived at the argument's image: in each, the previous lemma rewrites every
reached key to \code{Nat}'s block key, at which the subject's own clause applies.
\end{proof}

\begin{remark}
This reduces the capstone's per-argument environment conjunct at G8 to a single reachability fact,
\code{hargReach}, which stays a binder. That binder is \emph{not} slack: at the empty specification
environment the per-argument conjunct is false while the lowering of a bare constant still holds,
which is why the lemma carries the environment-erasure and lowering facts as antecedents rather than
being stated at every environment.
\end{remark}

\section{The eight rungs}
\label{sec:capstone:rungs}

\begin{center}
\begin{tabular}{llll}
rung & subject & answer & committed \code{.ast} \\
\hline
G1 & \code{spikeZero} & 0 & 354 bytes \\
G2 & \code{spikeLit} & 4 & 1556 bytes \\
G3 & \code{spikeLet} & 3 & 1469 bytes \\
G4 & \code{spikeProj} & 1 & 2284 bytes \\
G5 & \code{spikeCase} & 1 & 966 bytes \\
G6 & \code{spikeFix} & 2 & 853 bytes \\
G7 & \code{arithClosed} & 8 & 14497 bytes \\
G8 & \code{benchArith} applied to 0 & 8 & 14163 bytes \\
\end{tabular}
\end{center}

\begin{theorem}[Rung G1 is green]\label{thm:Green-green_G1}
\lean{LeanToLambdaBox.Green.green_G1} \leanok
\uses{thm:shipping_erase_correct_firstorder, def:Green-g1Table, def:Green-peanoLB, def:Witness-SourceTableAdequate, def:ErasureSpec, def:EraserAsks, def:UpstreamAsks, asm:htbl, asm:hsafe, asm:hblk, asm:hprep, asm:hrun, asm:hev, asm:hvwt, asm:hty, asm:hfo, asm:ErasureBridge-erasesEnv, asm:ErasureBridge-lowerEnv}
For the \code{\#erase} run recorded in \code{VerifyBench/Spikes/G1.lean} --- assuming the
specification bundle, an adequate and safe G1 table, this repository's and upstream's obligations,
the block condition, the preparation and run equations, the bridge, a source evaluation of the
subject to a value, that value's translation and typing at a \code{Nat} spine, and that \code{Nat}
is first order in the model --- there are a specification environment and an intermediate term such
that the subject erases to it, the environment erases the source environment along it, the lowering
carries both to the committed emitted program, that program is well formed for peregrine, and the
source value's erasure is unique, its lowering is the box-free literal \lbox{} numeral 0, and the
emitted program evaluates to that very numeral.
\end{theorem}

\begin{proof} \leanok
\uses{thm:shipping_erase_correct_firstorder, lem:Green-spike_configPinned, lem:Green-g1_compilerBodies, lem:Witness-trExprS_const_of_table, thm:supportedB_sound, lem:Green-g1_supported, lem:Green-g1_noBodylessRefs, lem:Green-g1_wf, lem:Green-g1_eval, thm:eval_deterministic, asm:lean4lean-trust, asm:lean-reflection-axioms}
The capstone is applied with six checked terms in place of six of its binders --- the pinned
configuration, G1's compiler bodies, the table route's translation, the fragment checker's verdict
transported by its soundness theorem, the body-less-reference check and the well-formedness check
--- and the remaining binders passed through. Its observable is then instantiated at the
\textbf{empty} argument list, where the length equation is \code{rfl}, the per-argument premise is
vacuous and the spine translation is again the table route. Finally the determinism of \lbox{}
evaluation identifies the existential target value with the committed literal answer, which is
substituted into the conclusion.
\end{proof}

\begin{remark}
G1 is the simplest rung and the only one at which \code{hcb} is a checked term; at every other rung
five checked terms are supplied and \code{hcb} is a binder.
\end{remark}

\begin{theorem}[Rungs G2, G3, G4 and G6]\label{thm:Green-green_G2}
\lean{LeanToLambdaBox.Green.green_G2, LeanToLambdaBox.Green.green_G3, LeanToLambdaBox.Green.green_G4, LeanToLambdaBox.Green.green_G6} \leanok
\uses{thm:shipping_erase_correct_firstorder, def:Green-g1Table, def:Green-peanoLB, def:Witness-SourceTableAdequate, def:CompilerBodies, asm:htbl, asm:hsafe, asm:hblk, asm:hcb, asm:hprep, asm:hrun, asm:hev, asm:hvwt, asm:hty, asm:hfo, asm:ErasureBridge-erasesEnv, asm:ErasureBridge-lowerEnv}
Four further rungs of the same shape as G1, at their own tables, subjects and committed programs,
with the literal answers 4, 3, 1 and 2, and with one extra binder each --- \code{hcb}, the
compiler-bodies hypothesis, which only G1 discharges. G2 exhibits a \code{Nat} literal becoming a
peano tower together with the \code{OfNat} class tower; G3 a \code{let} contracted by zeta on both
sides; G4 a projection taken on both sides with the pair's two type parameters erased to $\Box$; G6
a recursive constant whose emitted body is a fixpoint node that the target run unfolds twice under
the guard.
\end{theorem}

\begin{proof} \leanok
\uses{thm:shipping_erase_correct_firstorder, lem:Green-spike_configPinned, lem:Witness-trExprS_const_of_table, thm:supportedB_sound, lem:Green-g1_supported, lem:Green-g1_noBodylessRefs, lem:Green-g1_wf, lem:Green-g1_eval, thm:eval_deterministic, asm:lean4lean-trust, asm:lean-reflection-axioms}
Identical in shape to G1's: the capstone at the rung's five checked terms, the observable at the
empty spine, determinism to fix the answer. Together the four cover the erasure rules the fragment
admits --- literal, zeta, projection with boxed type parameters, delta on a fixpoint --- and at all
four both \code{hcb} and the source evaluation \code{hev} remain binders.
\end{proof}

\begin{theorem}[Rung G5 is green, and its source evaluation is not assumed]\label{thm:Green-green_G5}
\lean{LeanToLambdaBox.Green.green_G5} \leanok
\uses{thm:shipping_erase_correct_firstorder, def:Green-SpikeNatFacts, def:Green-SpikeUnitFacts, def:Green-peanoSrc, def:Green-peanoLB, def:Green-g1Table, def:StepDefeq, asm:htbl, asm:hsafe, asm:hblk, asm:hcb, asm:hprep, asm:hrun, asm:hvwt, asm:hty, asm:hfo, asm:ErasureBridge-erasesEnv, asm:ErasureBridge-lowerEnv}
As the other rungs, at G5's table, subject and committed program, except that the source-evaluation
binder is replaced by the two model-side bundles and the three definitional-step facts that the
derived source evaluation consumes, the observed value is fixed to the source numeral 1, and the
answer is the literal \lbox{} numeral 1.
\end{theorem}

\begin{proof} \leanok
\uses{thm:shipping_erase_correct_firstorder, prop:Green-g5_seval, lem:Green-spike_configPinned, lem:Witness-trExprS_const_of_table, thm:supportedB_sound, lem:Green-g1_supported, lem:Green-g1_noBodylessRefs, lem:Green-g1_wf, lem:Green-g1_eval, thm:eval_deterministic, asm:lean4lean-trust, asm:lean-reflection-axioms}
As at G1, with the capstone's source-evaluation premise supplied by
Proposition~\ref{prop:Green-g5_seval} rather than bound. What remains of the trust rows here is
\code{hcb}, the two value-side facts \code{hvwt} and \code{hty}, the first-order hypothesis, the two
model-side bundles and the three definitional-step facts.
\end{proof}

\begin{theorem}[Rung G7 is green]\label{thm:Green-green_G7}
\lean{LeanToLambdaBox.Green.green_G7} \leanok
\uses{thm:shipping_erase_correct_firstorder, def:Green-g1Table, def:Green-peanoLB, def:spikeZero, asm:htbl, asm:hsafe, asm:hblk, asm:hcb, asm:hprep, asm:hrun, asm:hev, asm:hvwt, asm:hty, asm:hfo, asm:ErasureBridge-erasesEnv, asm:ErasureBridge-lowerEnv}
As G1, at the tracked benchmark subject \code{arithClosed}, which is \code{benchArith 0}, with the
literal answer the \lbox{} peano numeral 8 and with \code{hcb} and \code{hev} as binders.
\end{theorem}

\begin{proof} \leanok
\uses{thm:shipping_erase_correct_firstorder, lem:Green-spike_configPinned, lem:Witness-trExprS_const_of_table, thm:supportedB_sound, lem:Green-g1_supported, lem:Green-g1_noBodylessRefs, lem:Green-g1_wf, lem:Green-g1_eval, thm:eval_deterministic, asm:lean4lean-trust, asm:lean-reflection-axioms}
As at G1, at G7's own checked terms. This is the ladder at real program scale --- forty emitted
declarations, ten class projections, four fixpoint blocks, five case nodes and a 14,497-byte
committed \code{.ast} --- and its well-formedness and body-less-reference
checks are the ladder's most expensive kernel computations.
\end{proof}

\begin{theorem}[Rung G8 is green, and it is the applied one]\label{thm:Green-green_G8}
\lean{LeanToLambdaBox.Green.green_G8} \leanok
\uses{thm:shipping_erase_correct_firstorder, def:Green-SpikeNatFacts, def:Green-peanoLB, def:Green-g1Table, def:ReachableFrom, def:ErasesEnv, def:LowerEnv, asm:htbl, asm:hsafe, asm:hblk, asm:hcb, asm:hprep, asm:hrun, asm:hargReach, asm:hev, asm:hvwt, asm:hty, asm:hfo, asm:ErasureBridge-erasesEnv, asm:ErasureBridge-lowerEnv}
As the others at G8's data, except that the subject \code{benchArith} is function-typed: the source
evaluation is read at the \textbf{spine} \code{benchArith Nat.zero}, the conclusion's target
evaluation at the emitted term \textbf{applied to} the \lbox{} numeral 0, and two binders appear
that the other seven do not have --- the \code{Nat} bundle, and \code{hargReach}, the hypothesis
that the subject's erasure reaches \code{Nat}'s block key in any specification environment carrying
the other facts. The answer is again the literal numeral 8.
\end{theorem}

\begin{proof} \leanok
\uses{thm:shipping_erase_correct_firstorder, lem:Green-g8_trExprS_spine, lem:Green-g8_argErasesEnv, lem:Green-spike_configPinned, lem:Witness-trExprS_const_of_table, thm:supportedB_sound, lem:Green-g1_supported, lem:Green-g1_noBodylessRefs, lem:Green-g1_wf, lem:Green-g1_eval, thm:eval_deterministic, def:Erases, def:Lower, asm:lean4lean-trust, asm:lean-reflection-axioms}
The capstone is applied as at the other rungs, then its observable is instantiated at the
\textbf{one-element} argument list, the source constant \code{Nat.zero} against the \lbox{} numeral
0. The per-argument premise is therefore not vacuous: it is discharged by the constructor arm of the
erasure relation and the constructor arm of the lowering relation at the numeral, together with
Lemma~\ref{lem:Green-g8_argErasesEnv}, whose reachability input is the binder \code{hargReach}. The
spine's translation is Lemma~\ref{lem:Green-g8_trExprS_spine}, and determinism fixes the answer to
the literal numeral 8. G8 is the only rung exercising a non-empty spine, hence the only one at which
the capstone's per-argument premise has content.
\end{proof}

\section{The tooling that re-measures}
\label{sec:capstone:tooling}

Nothing in this section is a graph node; all of it is machinery that re-derives, outside the kernel,
what the rungs assume or claim.

\code{lake exe green-check} re-runs the Lean frontend on a rung's spike module, byte-diffs the
regenerated \code{.ast} against the committed one \emph{and} against the literal program the rung's
theorem is stated about, and runs the certified evaluator on the rung's observed term against the
committed answer; \code{--self-test} exercises the evaluator on the delta, iota and fixpoint
fixtures independently of any rung. This is the external mechanisation of \code{hrun}.

\code{lake exe reify --check MOD NAME...} re-derives a committed table from the live environment
field by field --- the same comparison Definition~\ref{def:Witness-SourceTable-check} performs, run
outside the kernel --- tolerating bodies that agree only up to binder names and reporting those as
notes; \code{--blocks} re-checks the two computable clauses of \code{TableBlocks} against the
compiler's fix-block view; \code{--prepared} checks that preparation is the identity at a rung's
subject, which is the binder \code{hprep}. CI runs all three over the eight rung tables, and over a
deliberately stale table that must be \emph{rejected}.

\code{scripts/ledger.sh} runs \code{test/Ledger.lean} --- 40 \code{\#print axioms} invocations ---
and diffs the output against \code{test/ledger.expected}, so any change to a tracked theorem's
measured axiom footprint fails the build. \code{scripts/erases\_correct.sh} and
\code{scripts/erasesLB.sh} guard \emph{premise shape}, which a footprint cannot see.
\code{scripts/lean4lean-sorries.sh} re-scans the pinned lean4lean checkout and pins the inherited
trust boundary at sixteen \code{sorry} sites. \code{scripts/frozen.sh} holds the five benchmark
sources byte-identical against their expected copies and against the sibling benchmarks repository.

\code{lake exe coverage} regenerates \code{doc/coverage.md} by \emph{measuring} the tree ---
reifying the five programs' tables, running the fragment checker at every tabled body, reading the
eight rungs' data out of \code{Green.lean} --- and \code{--check} fails on drift. \code{lake exe
hygiene} enforces the documentation and dead-declaration policies, with a budget of 240 declarations
outside the import closure of \code{Green.lean} and \code{Capstone.lean}.

\section{Deviations from the paper, and caveats}
\label{sec:capstone:caveats}

The capstone's shape is [S, Sec. 7.3--7.4]'s, with three deliberate substitutions already named:
\lbox{}'s own weak call-by-value evaluation in place of MetaRocq's, the extra lowering relation
$\Lower$ with the consequent split of the environment side, and a source evaluation reading a
compiler-body table rather than the kernel's $\delta$ rules. Two further deviations are visible in
the emitted data: constructors are emitted in \emph{applied} form, never as blocks (F-ETA2), and no
fixpoint eta-expansion is performed, so the emitted program satisfies \code{LBWfPeregrine} but not
the stronger expanded-fixpoint condition (F-ETA).

The honest reading of the ladder is this. Eight rungs instantiate the capstone at real
\code{\#erase} runs and end in literal answers, so the theorem is not vacuously satisfiable by
$\Box$ or by a stuck term, and every hypothesis a computation can settle --- the configuration, the
fragment check, the body-less-reference check, the well-formedness check, the target evaluation, and
at G1 the compiler bodies, at G5 the source evaluation --- is settled by one. But \code{lenv} and
\code{env} are universally quantified at every rung; the run equation \code{hrun} is an equation
about an opaque world token that no Lean term inhabits; the binder \code{hbridge} is inhabited
nowhere; \code{hcb} is a binder at seven rungs and \code{hev} at seven; and the \code{PUnit} half of
G5's model-side bundle is satisfied by no fixture. Two fields of \code{EraserAsks}, a binder of the
capstone and of every rung, are moreover refuted in general: the kernel inductive-head field is
false at a telescope of at least 256 binders (finding F-DEPTH, since Lean's approximate depth field
is eight bits wide), and the block-key-distinctness field is false on a legal mutual \code{unsafe
def} block whose \code{\_unsafe\_rec} companion collides (finding F-UNSAFEREC). Both are true on
every measured input, and both bound what "verified" may mean here. Chapter~\ref{chap:trust} states
each of these as its own assumption and gives the measured axiom footprint behind them.
